% Daten – bank/daten/ (zusammenbau v0.4, Heft)
\einheitenkopf[t1]{Daten}
\setcounter{aufgabe}{0}

% Anlauf (ohne Prüfkennung)
\begin{aufgabe}{Häufigkeiten – Anlauf}
\begin{teile}
\teil Von $26$ Kindern fahren $8$ mit dem Bus. – Was ist das Ganze? Unterstreiche es.
\teil Zähle die Strichliste aus: Wie oft wurde Rot genannt, wie oft Blau?

\sachtabelle{lc}{Farbe & Strichliste}{Rot & \strichliste{11} \\ Blau & \strichliste{7}}
\teil Würfe: $3$, $5$, $2$, $5$, $1$, $5$, $3$, $2$, $5$, $3$. Trage die Anzahl je Augenzahl in die Tabelle ein.

\sachtabelle{lcccc}{Augenzahl & 1 & 2 & 3 & 5}{Anzahl & \leerzelle & \leerzelle & \leerzelle & \leerzelle}
\end{teile}
\end{aufgabe}

\begin{aufgabe}{Häufigkeiten – Prüfungsaufgaben}
\begin{teile}
\teil Eine Handballmannschaft spielte in einer Saison $26$ Spiele: $17$ gewonnen, $6$ unentschieden, $3$ verloren. Gib die relative Häufigkeit eines Siegs als Bruch und in Prozent an (auf eine Stelle gerundet). \hfill (P10 2015 OS)
\teil Eine Volleyballmannschaft spielte $28$ Spiele: $16$ gewonnen, $7$ unentschieden, $5$ verloren. Gib die relative Häufigkeit eines Siegs als Bruch und in Prozent an (auf eine Stelle gerundet). \hfill (P10 2015 OS)
\end{teile}
\end{aufgabe}

% Anlauf (ohne Prüfkennung)
\begin{aufgabe}{Diagramme lesen – Anlauf}
\begin{teile}
\teil Zwischen zwei beschrifteten Werten der Achse liegen mehrere Hilfslinien. Was ist eine Hilfslinie wert? Gib nur den Schritt an.

\saeulenab[ymax=600,ystep=100,yfein=20,ylabel=Anzahl]{A/280,B/460,C/140}
\leerfeld je Linie
\teil Lies ab: Wie viele Besucher kamen am Dienstag?

\saeulenab[ymax=80,ystep=10,yfein=10,ylabel=Besucher]{Mo/30,Di/40,Mi/70,Do/50}
\leerfeld Besucher
\teil Lies ab: Wie viele Punkte hat Team B?

\saeulenab[ymax=100,ystep=20,yfein=5,ylabel=Punkte]{A/40,B/65,C/90}
\leerfeld Punkte
\end{teile}
\end{aufgabe}

\begin{aufgabe}{Diagramme lesen – Prüfungsaufgaben}
\begin{teile}
\teil Lies ab: Wie viele Besucher hatte das Museum im Jahr 2014? Gib die Zahl vollständig an. \hfill (P10 2016 OS) \\

\saeulenab[ymax=400,ystep=100,yfein=10,ylabel=Besucher in Tausend]{2012/310,2013/240,2014/270,2015/330}
\leerfeld Besucher
\teil Lies ab: An welchen Tagen kamen mehr als $220$ Gäste ins Freibad? \hfill (P10 2020 OS)

\balkenab[xmax=320,xstep=40,xfein=10,xlabel=Gäste]{Mo/160,Di/240,Mi/110,Do/190,Fr/260,Sa/220,So/300}
\teil Lies ab: In welchen Jahren hatte der Zoo weniger als $400\,000$ Besucher? \hfill (P10 2016 OS)

\saeulenab[ymax=600,ystep=100,yfein=20,ylabel=Besucher in Tausend]{2018/440,2019/380,2020/520,2021/360,2022/400,2023/480}
\teil Zeichne den Balken für Mittwoch: $230$ Gäste. \hfill (P10 2020 OS)

\balkenab[xmax=300,xstep=50,xfein=10,xlabel=Gäste]{Mo/160,Di/260,Mi/,Do/220,Fr/110}
\teil Zeichne die Säule für Sonntag: $1{,}45$ mm Regen. \hfill (P10 2026 FOR)

\saeulenab[ymax=2.5,ystep=0.5,yfein=0.05,ylabel=Regen in mm]{Mo/1.2,Di/1.7,Mi/0.9,Do/2.3,Fr/1.6,Sa/1.1,So/}
\end{teile}
\end{aufgabe}

\begin{aufgabe}{Diagramme lesen – Prüfungsaufgaben – weiter}
\begin{teile}
\teil Einwohner in Millionen: Polen $38$, Spanien $47$, Italien $59$, Frankreich $68$, Deutschland $84$. Die y-Achse trägt keine Zahlen; die Säule für Polen ist $3{,}8$ Kästchen hoch. a) Wie viel Mio. ist ein Kästchen wert? Beschrifte die Achse. b) Welches Land zeigt die Säule mit dem Fragezeichen ($5{,}9$ Kästchen)? c) Zeichne die Säule für Frankreich. \hfill (P10 2023 OS)

\saeulenab[ymax=90,ystep=10,yfein=10,ylabel=Einwohner in Mio.,ohnezahlen]{Polen/38,?/59,Frankreich/,Deutschland/84}
% AUSGELASSEN daten-e2-k2-s11-v2: Dimension too large.
% \teil Sitzplätze in Stadien: A $81\,000$, B $75\,000$, C $60\,000$, D $54\,000$, E $42\,000$. Die y-Achse trägt keine Zahlen; die Säule für D ist $5{,}4$ Kästchen hoch. a) Wie viele Plätze ist ein Kästchen wert? Beschrifte die Achse. b) Welches Stadion zeigt die Säule mit dem Fragezeichen ($7{,}5$ Kästchen)? c) Zeichne die Säule für E. \hfill (P10 2023 OS)
% 
% \saeulenab[ymax=90000,ystep=10000,yfein=10000,ylabel=Sitzplätze,ohnezahlen]{A/81000,?/75000,D/54000,E/}
\teil \textit{(ausgelassen)}
\teil Tabelle der verkauften Früchte: Apfel $84$, Banane $48$, Kirsche $72$, Birne $84$, Traube $36$. Das Diagramm zeigt fünf Säulen A bis E ohne Namen, die Achse trägt nur Gitterlinien; die längsten Säulen B und E sind $7$ Linien hoch. a) Wie viel ist eine Gitterlinie wert? b) Ordne die Säulen den Früchten zu. \hfill (P10 2018 OS)

\saeulenab[ymax=96,ystep=12,yfein=12,ylabel=Anzahl,ohnezahlen]{A/72,B/84,C/48,D/36,E/84}
\teil Tabelle der Haustiere in einer Schule: Hund $150$, Katze $180$, Hamster $60$, Fisch $90$, Vogel $60$. Das Diagramm zeigt fünf Säulen A bis E ohne Namen, die Achse trägt nur Gitterlinien; die längste Säule C ist $6$ Linien hoch. a) Wie viel ist eine Gitterlinie wert? b) Ordne die Säulen den Tieren zu. \hfill (P10 2018 OS)

\saeulenab[ymax=210,ystep=30,yfein=30,ylabel=Anzahl,ohnezahlen]{A/90,B/60,C/180,D/60,E/150}
\end{teile}
\end{aufgabe}

% Anlauf (ohne Prüfkennung)
\begin{aufgabe}{Anteile darstellen – Anlauf}
\begin{teile}
\teil Zu einem Kreisdiagramm heißt es: „Der Sektor Bus hat einen Mittelpunktswinkel von $54$.“ Was ist die Zahl $54$ hier? Kreuze an.\\ \kreuz{Anteil in \%} \\ \kreuz{Winkel in Grad}
\teil Färbe $30\,\%$ im Streifen. Wie lang ist der Abschnitt?

\streifenleer
\leerfeld[cm]
\teil Zeichne die Anteile nacheinander in den Streifen und beschrifte sie: Rad $40\,\%$, Bus $30\,\%$, zu Fuß $30\,\%$.

\streifenleer[0]
\end{teile}
\end{aufgabe}

\begin{aufgabe}{Anteile darstellen – Prüfungsaufgaben}
\begin{teile}
\teil Zeichne die Anteile nacheinander in den Streifen und beschrifte sie: Sport $20\,\%$, Musik $45\,\%$, Lesen $10\,\%$, Spiele $25\,\%$. \hfill (P10 2018 OS)

\streifenleer[0]
\teil Frachtanteile eines Hafens: Kohle $21\,\%$, Erz $58\,\%$, Holz $7\,\%$, Sonstiges $2\,\%$, Getreide ? – zusammen $100\,\%$. Wie groß ist der Anteil Getreide? \hfill (P10 2021 OS) \\ \leerfeld[\%]
\teil Von $76$ Elfmetern wurden $9$ gehalten. Welcher Mittelpunktswinkel gehört im Kreisdiagramm zu „gehalten“? Runde auf ganze Grad. \hfill (P10 2015 OS) \\ \leerfeld °
\teil $640$ Befragte lesen Zeitung: täglich $42\,\%$, mehrmals pro Woche $23\,\%$, selten $27\,\%$, nie $8\,\%$. Nur „nie“ ist beschriftet. Beschrifte die drei anderen Sektoren und berechne den Mittelpunktswinkel für „nie“. \hfill (P10 2022 OS)

\kreisdiagramm{/42,/27,/23,nie/8}
\end{teile}
\end{aufgabe}

\begin{aufgabe}{Anteile darstellen – Prüfungsaufgaben – weiter}
\begin{teile}
\steil Stromverbrauch eines Haushalts in einem Jahr: $3\,600$ kWh, davon Heizung $1\,440$, Warmwasser $720$, Kühlen $540$, Kochen $360$, Sonstiges $360$, Licht $180$ (kWh). Das Kreisdiagramm zeigt die sechs Anteile; nur Heizung und Licht sind beschriftet. a) Beschrifte den Sektor „Kühlen“. b) Berechne den Mittelpunktswinkel für „Licht“. \hfill (P10 2024 OS)

\kreisdiagramm{Heizung/1440,/720,/540,/360,/360,Licht/180}
\steil Wasserverbrauch einer Familie an einem Tag: $420$ Liter, davon Dusche $168$, Toilette $105$, Wäsche $63$, Küche $42$, Sonstiges $28$, Garten $14$ (Liter). Das Kreisdiagramm zeigt die sechs Anteile; nur Dusche und Garten sind beschriftet. a) Beschrifte den Sektor „Wäsche“. b) Berechne den Mittelpunktswinkel für „Garten“. \hfill (P10 2024 OS)

\kreisdiagramm{Dusche/168,/105,/63,/42,/28,Garten/14}
\teil Anteile der Verkehrsmittel: Auto $52{,}6\,\%$, Bahn $29{,}8\,\%$, Rad $12{,}4\,\%$, zu Fuß $5{,}2\,\%$. Im Kreisdiagramm ist erst ein kleiner Sektor von etwa $19^\circ$ eingezeichnet. a) Welches Verkehrsmittel gehört zu diesem Sektor? b) Berechne den Mittelpunktswinkel für „Rad“ (ganze Grad). c) Zeichne den Sektor „Rad“ an den vorhandenen an und beschrifte ihn. \hfill (P10 2017 OS)

\kreisdiagramm{/19,/341}
\teil Anteile der Energieträger: Wind $41{,}7\,\%$, Sonne $33{,}5\,\%$, Wasser $18{,}3\,\%$, Biomasse $6{,}5\,\%$. Im Kreisdiagramm ist erst ein kleiner Sektor von etwa $23^\circ$ eingezeichnet. a) Welcher Energieträger gehört zu diesem Sektor? b) Berechne den Mittelpunktswinkel für „Wasser“ (ganze Grad). c) Zeichne den Sektor „Wasser“ an den vorhandenen an und beschrifte ihn. \hfill (P10 2017 OS)

\kreisdiagramm{/23,/337}
\teil Alter der $13$ Trainerinnen und Trainer eines Vereins: $27$, $34$, $41$, $29$, $52$, $38$, $45$, $31$, $60$, $36$, $48$, $55$, $43$. a) Wie viel Prozent sind jünger als $40$ Jahre? Runde auf eine Stelle. b) Zeichne diesen Anteil als Sektor in das Kreisdiagramm ein. \hfill (P10 2025 OS)

\kreisleer
\teil Körpergrößen der $9$ Spieler einer Mannschaft in cm: $178$, $185$, $172$, $190$, $181$, $176$, $188$, $179$, $183$. a) Wie viel Prozent sind größer als $180$ cm? Runde auf eine Stelle. b) Zeichne diesen Anteil als Sektor in das Kreisdiagramm ein. \hfill (P10 2025 OS)

\kreisleer
\end{teile}
\end{aufgabe}

% Anlauf (ohne Prüfkennung)
\begin{aufgabe}{Kenngrößen – Anlauf}
\begin{teile}
\teil „Wie weit liegen der kleinste und der größte Wert auseinander?“ – Welche Kenngröße ist gefragt? Kreuze an.\\ \kreuz{Spannweite} \\ \kreuz{Modalwert} \\ \kreuz{Median} \\ \kreuz{Mittelwert}
\teil Punkte: $14$, $9$, $22$, $17$, $11$. Minimum und Maximum? Min = \leerfeld , Max = \leerfeld
\teil Der Durchschnitt von fünf Werten ist $20$. Ein sechster Wert $32$ kommt dazu. Steigt oder sinkt der Mittelwert? Begründe ohne Neuberechnung.
\end{teile}
\end{aufgabe}

\begin{aufgabe}{Kenngrößen – Prüfungsaufgaben}
\begin{teile}
\teil Lies aus der Tabelle die kleinste und die größte Wochenzahl für Super E10 ab (nicht die Dieselzeile). \hfill (P10 2014 OS) \\

\sachtabelle{lcccccc}{Woche & 1 & 2 & 3 & 4 & 5 & 6}{Super E10 in ct & 168{,}3 & 171{,}9 & 165{,}0 & 162{,}4 & 169{,}7 & 166{,}2 \\ Diesel in ct & 155{,}1 & 158{,}6 & 151{,}9 & 149{,}8 & 156{,}3 & 153{,}0}
Min = \leerfeld , Max = \leerfeld
\teil Muttersprachler in Millionen: Portugiesisch $230$, Bengalisch $240$, Russisch $150$, Japanisch $125$, Deutsch $95$. Gib Minimum, Maximum und Spannweite an. \hfill (P10 2023 OS) \\ Min = \leerfeld , Max = \leerfeld , Spannweite = \leerfeld
\teil Weitsprung in m: $4{,}35$; $4{,}62$; $4{,}18$; $4{,}50$; $4{,}27$; $4{,}44$. Spannweite? \hfill (P10 2019 OS) \\ \leerfeld[m]
\teil Lies aus der Tabelle die Spannweite des Dieselpreises ab (Dieselzeile). \hfill (P10 2014 OS) \\

\sachtabelle{lcccccc}{Woche & 1 & 2 & 3 & 4 & 5 & 6}{Super E10 in ct & 168{,}3 & 171{,}9 & 165{,}0 & 162{,}4 & 169{,}7 & 166{,}2 \\ Diesel in ct & 155{,}1 & 158{,}6 & 151{,}9 & 149{,}8 & 156{,}3 & 153{,}0}
\leerfeld ct
\teil Dieselpreise in € je Liter: Mo $1{,}62$; Di $1{,}64$; Mi $1{,}61$; Do $1{,}66$; Fr $1{,}71$; Sa $1{,}69$; So $1{,}73$. Spannweite? \hfill (P10 2026 FOR) \\ \leerfeld[€]
\teil Messwerte in °C: $6$; $9$; $2$; $11$; $4$; $9$; $7$. Bestimme den Zentralwert (Median). \hfill (P10 2015 OS)
\end{teile}
\end{aufgabe}

\begin{aufgabe}{Kenngrößen – Prüfungsaufgaben – weiter}
\begin{teile}
\teil Datenreihe in °C: $16$, $13$, $19$, $14$, $19$, $9$. Bestimme den Median. \hfill (P10 2024 OS)
\teil Werte: $6$; $30$; $45$. Arithmetisches Mittel? \hfill (P10 2016 OS)
\teil Sprunghöhen: $1{,}45$ m; $1{,}50$ m; $1{,}35$ m; $1{,}50$ m. Arithmetisches Mittel? \hfill (P10 2017 OS) \\ \leerfeld[m]
\teil Weiten: $5{,}12$ m; $4{,}96$ m; $5{,}04$ m; $5{,}28$ m. Durchschnittliche Weite? \hfill (P10 2021 OS) \\ \leerfeld[m]
\teil Bewegte Container eines Hafens in Mio. 2016 bis 2023: $6{,}2$; $6{,}9$; $7{,}4$; $7{,}1$; $7{,}8$; $8{,}3$; $7{,}9$; $8{,}0$. Berechne Spannweite und arithmetisches Mittel (zwei Stellen). \hfill (P10 2021 OS) \\ \leerfeld Mio., \leerfeld Mio.
\end{teile}
\end{aufgabe}

\begin{aufgabe}{Kenngrößen – Prüfungsaufgaben – weiter}
\begin{teile}
\teil Lies die Besucherzahlen ab; der Wert für Donnerstag ($150$) steht nur im Text. Durchschnittliche Besucherzahl je Tag? \hfill (P10 2020 OS)

\saeulenab[ymax=300,ystep=50,yfein=10,ylabel=Besucher]{Mo/120,Di/90,Mi/160,Do/,Fr/140,Sa/210,So/250}
\teil Jährlich werden $900$ Jugendliche befragt, wie viele täglich Sport treiben. Lies ab: Minimum und Maximum der fünf Jahre, dann das arithmetische Mittel. \hfill (P10 2022 OS)

\saeulenab[ymax=400,ystep=100,yfein=10,ylabel=Jugendliche]{2020/350,2021/390,2022/340,2023/310,2024/360}
\teil Monatliche Niederschläge eines Jahres in mm. Lies die Werte ab und berechne die Spannweite. \hfill (P10 2024 OS)

\saeulenab[ymax=100,ystep=20,yfein=2,ylabel=Niederschlag in mm]{J/40,F/16,M/34,A/28,M/56,J/72,J/92,A/68,S/44,O/50,N/62,D/38}
\teil Zu $19$ Heimspielen kamen insgesamt $412\,879$ Zuschauer. Durchschnittliche Zuschauerzahl je Spiel, sinnvoll gerundet? \hfill (P10 2015 OS)
\teil In einem Jahr fielen insgesamt $634{,}8$ mm Niederschlag. Behauptung: „Im Monatsdurchschnitt sind das etwa $52{,}9$ mm.“ Weise die Behauptung rechnerisch nach. \hfill (P10 2024 OS)
\teil Tageswerte in °C: Mo $17$, Di $19$, Mi $16$, Do $21$, Fr $18$, Sa ?, So $15$. Der Wochendurchschnitt beträgt $18$ °C. Trage den fehlenden Samstagswert ein. \hfill (P10 2022 OS) \\ Sa: \leerfeld °C
\end{teile}
\end{aufgabe}

\begin{aufgabe}{Kenngrößen – Prüfungsaufgaben – weiter}
\begin{teile}
\teil Alter der Kinder einer Gruppe: $6$, $3$, $2$, $5$, $3$, $5$, $3$, $5$, $3$, $4$. Aussage 1: „Die Spannweite beträgt $4$ Jahre.“ Aussage 2: „Der Modalwert beträgt $5$ Jahre.“ Welche Aussage ist falsch? Kreuze an und berichtige sie. \hfill (P10 2025 OS)\\ \kreuz{Aussage 1} \\ \kreuz{Aussage 2}
\teil Besucher je Tag: Mo $120$, Di $90$, Mi $160$, Do $150$, Fr $140$, Sa $210$, So $250$. Aussage 1: „Die Spannweite beträgt $160$.“ Aussage 2: „Am Sonntag kamen ein Viertel mehr Besucher als am Mittwoch.“ Prüfe beide Aussagen rechnerisch. \hfill (P10 2020 OS)
\teil Dieselpreise in € je Liter: Mo $1{,}62$; Di $1{,}64$; Mi $1{,}61$; Do $1{,}66$; Fr $1{,}71$; Sa $1{,}69$; So $1{,}73$. Behauptung: „Von Montag bis Freitag liegt der Durchschnittspreis unter $1{,}65$ €.“ Prüfe die Behauptung rechnerisch. \hfill (P10 2026 FOR)
\steil Alter der $11$ Mitglieder einer Theatergruppe: $47$, $23$, $39$, $61$, $35$, $29$, $52$, $44$, $19$, $58$, $33$. a) Berechne das Durchschnittsalter. b) Die älteste und die jüngste Person verlassen die Gruppe. Wie wirkt sich das auf das Durchschnittsalter aus? Begründe. \hfill (P10 2025 OS)
\steil Alter der $11$ Spielerinnen einer Mannschaft: $24$, $31$, $19$, $27$, $35$, $22$, $29$, $26$, $33$, $21$, $30$. a) Berechne das Durchschnittsalter. b) Die älteste und die jüngste Spielerin wechseln den Verein. Wie wirkt sich das auf das Durchschnittsalter aus? Begründe. \hfill (P10 2025 OS)
\end{teile}
\end{aufgabe}

% Anlauf (ohne Prüfkennung)
\begin{aufgabe}{Beurteilen – Anlauf}
\begin{teile}
\teil Die y-Achse eines Säulendiagramms ist mit $0$, $50$, $100$, $150$ beschriftet. Wo beginnt sie? Kreuze an.\\ \kreuz{bei null} \\ \kreuz{nicht bei null}
\teil Aussage: „Im März wurden mehr als $60$ Bücher ausgeliehen.“ Stimmt das? Prüfe am Diagramm.

\saeulenab[ymax=100,ystep=20,yfein=5,ylabel=Bücher]{Jan/45,Feb/75,Mär/60,Apr/85}
% AUSGELASSEN daten-e5-k1-s6-v1: Dimension too large.
% \teil Ein Sparbuch bringt $0{,}5\,\%$ Zinsen im Jahr. Das Diagramm zeigt das Guthaben von 2020 bis 2024; es scheint stark zu wachsen. Erkläre, wie das Diagramm täuscht.
% 
% \begin{ksys}[xmin=2019,xmax=2025,ymin=1990,ymax=2050,xstep=1,ystep=10,xlabel=Jahr,ylabel=Guthaben in €,ablesen] \punkt{2020}{2000}{} \punkt{2021}{2010}{} \punkt{2022}{2020}{} \punkt{2023}{2030}{} \punkt{2024}{2040}{} \end{ksys}
\teil \textit{(ausgelassen)}
\end{teile}
\end{aufgabe}

\begin{aufgabe}{Beurteilen – Prüfungsaufgaben}
\begin{teile}
\teil Jährlich werden $900$ Jugendliche befragt, wie viele täglich Sport treiben. Behauptung: „Von Jahr zu Jahr treiben immer weniger Jugendliche täglich Sport. Von 2020 bis 2023 hat sich die Anzahl um mehr als die Hälfte verringert.“ Prüfe beide Teile und begründe. \hfill (P10 2022 OS)

\saeulenab[ymax=400,ystep=100,yfein=10,ylabel=Jugendliche]{2019/372,2020/348,2021/361,2022/330,2023/296}
\end{teile}
\end{aufgabe}

\begin{aufgabe}{Beurteilen – Prüfungsaufgaben – weiter}
\begin{teile}
% AUSGELASSEN daten-e5-k1-s10-v1: Dimension too large.
% \steil Ein Artikel (2023) schreibt: „Die Zahl der Jugendlichen, die ein Instrument spielen, nimmt von Jahr zu Jahr dramatisch ab. In wenigen Jahren spielt kein Jugendlicher mehr ein Instrument.“ Das Diagramm zeigt die Anteile von 2019 bis 2023. Nenne eine der unzutreffenden Aussagen und berichtige sie. \hfill (P10 2019 OS)
% 
% \begin{ksys}[xmin=2018,xmax=2024,ymin=25,ymax=35,xstep=1,ystep=1,xlabel=Jahr,ylabel=Anteil in \%,ablesen] \punkt{2019}{34}{} \punkt{2020}{33}{} \punkt{2021}{31}{} \punkt{2022}{30}{} \punkt{2023}{28}{} \end{ksys}
\steil \textit{(ausgelassen)}
% AUSGELASSEN daten-e5-k1-s10-v2: Dimension too large.
% \steil Ein Artikel (2024) schreibt: „Immer weniger Menschen fahren mit der Bahn – der Anteil bricht ein. Bald wird niemand mehr Bahn fahren.“ Das Diagramm zeigt die Anteile von 2020 bis 2024. Nenne eine der unzutreffenden Aussagen und berichtige sie. \hfill (P10 2019 OS)
% 
% \begin{ksys}[xmin=2019,xmax=2025,ymin=18,ymax=24,xstep=1,ystep=1,xlabel=Jahr,ylabel=Anteil in \%,ablesen] \punkt{2020}{22}{} \punkt{2021}{22}{} \punkt{2022}{21}{} \punkt{2023}{20}{} \punkt{2024}{19}{} \end{ksys}
\steil \textit{(ausgelassen)}
\teil Lena zeichnet die Dieselpreise der Woche in ein Säulendiagramm, dessen y-Achse bei $2{,}10$ € beginnt. Sie sagt: „Das Diagramm zeigt, dass sich der Preis von Donnerstag ($2{,}14$ €) auf Freitag ($2{,}18$ €) verdoppelt.“ Erkläre, wie dieser falsche Eindruck entsteht. \hfill (P10 2026 FOR)
\teil Ein Diagramm zeigt die Temperatur an vier Tagen; die y-Achse beginnt bei $18$ °C. Jemand sagt: „Am Mittwoch ($21$ °C) war es dreimal so warm wie am Montag ($19$ °C), das sieht man an den Säulen.“ Erkläre, wie dieser falsche Eindruck entsteht. \hfill (P10 2026 FOR)
\steil Umfrage unter Frauen, wie oft sie in der Woche Sport treiben: nie $4\,\%$, 1- bis 2-mal $28\,\%$, 3- bis 4-mal $52\,\%$, öfter $16\,\%$. Aussage 1: „Weniger als die Hälfte der Frauen treiben 3- bis 4-mal Sport.“ Aussage 2: „Jede 25. Frau treibt nie Sport.“ Kreuze je Aussage an und begründe. \hfill (P10 2018 OS)\\ Aussage 1: \kreuz{richtig} \kreuz{falsch} \kreuz{nicht entscheidbar} \\ Aussage 2: \kreuz{richtig} \kreuz{falsch} \kreuz{nicht entscheidbar}
\steil Umfrage unter Jugendlichen, wie oft sie in der Woche kochen: nie $10\,\%$, 1- bis 2-mal $45\,\%$, 3- bis 4-mal $30\,\%$, öfter $15\,\%$. Aussage 1: „Mehr als die Hälfte kochen höchstens 2-mal in der Woche.“ Aussage 2: „Jeder 10. Jugendliche kocht 3- bis 4-mal.“ Kreuze je Aussage an und begründe. \hfill (P10 2018 OS)\\ Aussage 1: \kreuz{richtig} \kreuz{falsch} \kreuz{nicht entscheidbar} \\ Aussage 2: \kreuz{richtig} \kreuz{falsch} \kreuz{nicht entscheidbar}
\end{teile}
\end{aufgabe}
